\documentclass[11pt,a4paper]{article}
\usepackage{fontspec,amsmath,amssymb,amsthm,microtype}
\setmainfont{lmroman10-regular.otf}[BoldFont=lmroman10-bold.otf,ItalicFont=lmroman10-italic.otf,BoldItalicFont=lmroman10-bolditalic.otf]
\usepackage[margin=25mm]{geometry}
\usepackage[unicode,hidelinks]{hyperref}
\hypersetup{pdfauthor={LangevinLean Project}}
\setlength{\emergencystretch}{3em}
\setlength{\parskip}{0.4em}
\setlength{\parindent}{0pt}
\allowdisplaybreaks[2]
\raggedbottom
\title{\texorpdfstring{Construction of Gaussian coordinates with Brownian covariance}{Construction of Gaussian coordinates with Brownian covariance}}
\author{LangevinLean Project}
\date{}
\begin{document}\maketitle
Draft mathematical exposition; these notes may contain mistakes.


In the time Hilbert space \(L^2(\mathbb R,dt)\), define \(e_t=1_{(0,t]}\) for \(t\ge0\). Then \(\langle e_s,e_t\rangle=\min(s,t)\). Applying the isonormal map gives centered Gaussian variables \(B_t\) with this covariance. Linearity gives joint Gaussianity, and disjoint increment indicators are orthogonal, so the corresponding Gaussian increments are independent.



Independent product copies provide the spatial coordinates. Fixed-time representatives are measurable and retain the stated Gaussian laws. The conclusions concern finite-dimensional distributions and increments; sample-path continuity is not asserted. The later integral equation uses continuity in the mean-square norm, proved from the increment moments.


\section{\texorpdfstring{The time-\allowbreak{}function Hilbert space}{The time-function Hilbert space}}


Take the deterministic Hilbert space \(L^2(\mathbb R,dt)\). Interval indicators are its vectors, and the isonormal covariance formula evaluates their inner products. It is distinct from the random-variable \(L^2\) space on the Gaussian source.



The time-index Hilbert space uses Lebesgue square integration on the real line. Indicators of finite intervals belong to it because their squared integral is the interval length.


\section{\texorpdfstring{An interval indicator as a time index}{An interval indicator as a time index}}


Let \(e_t=\mathbf1_{(0,t]}\) as a class in \(L^2(dt)\), for \(t\ge0\). Construct it with the square-integrable interval indicator and the finite Lebesgue measure of \((0,t]\). Endpoints do not affect this Lebesgue \(L^2\) class, but this half-open convention makes interval intersection identities direct.



The indicator is one before time \(t\) and zero afterward, up to sets of Lebesgue measure zero. Its squared Hilbert norm is \(t\), so the associated isonormal variable has variance \(t\).


\section{\texorpdfstring{The overlap inner product of two indicators}{The overlap inner product of two indicators}}


The product of the two indicators is the indicator of \([0,\min(s,t)]\) almost everywhere. Integrating it gives \(\langle e_s,e_t\rangle=\min(s,t)\). The proof evaluates the Lebesgue measure of that finite interval, giving the deterministic inner product used in the Gaussian covariance calculation.



The product of the two indicators is the indicator of their overlap. Its integral is \(\min(s,t)\), yielding the desired covariance kernel exactly.


\section{\texorpdfstring{A countable basis index}{A countable basis index}}


Choose an orthonormal basis of the time-function Hilbert space and use its index set to label the independent scalar Gaussian coordinates. Orthogonal-series synthesis depends only on square-summable coefficients, so its definition does not require fixing an enumeration.



A countable basis index permits synthesis from a countable Gaussian product source. The particular enumeration affects the construction's description but not its covariance or resulting Gaussian laws.


\section{\texorpdfstring{An orthonormal basis for time functions}{An orthonormal basis for time functions}}


Select a Hilbert basis for the Hilbert space of time functions by Hilbert-space basis existence. This basis supplies the coordinate isometry for the explicit isonormal Gaussian family; all subsequent covariance identities follow from isometry and do not depend on which basis was selected.



An orthonormal basis expands each time indicator with square-summable coefficients. Parseval identifies their squared sum with the interval length, ensuring the Brownian synthesis has the correct variance.


\section{\texorpdfstring{The scalar Gaussian sample space}{The scalar Gaussian sample space}}


A sample is a real family indexed by the chosen time-Hilbert basis. Equip this product space with the independent standard Gaussian coordinate measure used in the isonormal series.



The sample space is the Gaussian coordinate source underlying the isonormal map. Time is an index of synthesized variables on this fixed space, rather than an additional random coordinate.


\section{\texorpdfstring{The scalar source probability}{The scalar source probability}}


Take the product of standard scalar Gaussian measures indexed by the chosen time basis. This probability has independent coordinate projections, providing the source variables used to synthesize every time value on one common sample space.



The scalar source probability has mass one. All synthesized time values are therefore random variables under one probability law, permitting their finite-dimensional covariances to be compared.


\section{\texorpdfstring{Normalization of the scalar probability}{Normalization of the scalar probability}}


The scalar source is the Gaussian coordinate product probability. Reuse normalization of the independent Gaussian product probability to provide normalization for the chosen basis index. This normalization enables the \(L^2\) expectation and covariance theorems for the constructed Brownian process.



Normalization of the source is inherited from its Gaussian product construction. This supplies the common probability assumption for Gaussian process and moment statements that follow.


\section{\texorpdfstring{A time value in the square-\allowbreak{}integrable space}{A time value in the square-integrable space}}


Apply the isonormal isometry to \(e_t\), obtaining \(B_t=I(e_t)\) as an \(L^2\) random variable. The construction immediately gives a centered Gaussian with variance \(\|e_t\|^2=t\); joint Gaussianity follows from linearity.



Apply the isonormal map to the interval indicator to obtain a scalar \(L^2\) time value. Its centered Gaussian law and variance follow from the indicator norm calculation.


\section{\texorpdfstring{The scalar Gaussian process}{The scalar Gaussian process}}


At each time select the strongly measurable representative of a brownian value as a square-integrable variable. This defines a measurable random variable at each time, written \(B_t(\omega)\) for use in the stated pushforward law and finite-dimensional distributions. Equalities of \(L^2\) classes still translate to almost-everywhere equalities of these representatives.



Choose the prescribed representative of each \(L^2\) time value to define the scalar process. Fixed-time laws are unaffected by the representative choice on null sets.


\section{\texorpdfstring{Fixed-\allowbreak{}time measurability of the process}{Fixed-time measurability of the process}}


The representative selected for an \(L^2\) class is measurable in this setting. Apply the representative measurability theorem to the square-integrable value at time \(t\), obtaining measurability of the time-\(t\) random variable.



At each fixed time the selected representative is measurable. This is the measurability needed for its marginal pushforward law and for assembling finite collections of time values.


\section{\texorpdfstring{The mean of the process}{The mean of the process}}


The isonormal mean theorem gives \(\mathbb EI(e_t)=0\). Since the process value at time \(t\) is precisely the selected representative of this class, this is the expectation of the displayed random function \(B_t\).



The isonormal map is centered at every index vector. Applying it to a time indicator therefore gives zero mean for each process value.


\section{\texorpdfstring{The Brownian covariance formula}{The Brownian covariance formula}}


The isonormal covariance formula gives \(\operatorname{Cov}(I(e_s),I(e_t))=\langle e_s,e_t\rangle\). Substitute the overlap formula \(\langle1_{(0,s]},1_{(0,t]}\rangle=\min(s,t)\) to obtain \(\min(s,t)\). This calculation establishes the Brownian covariance for the chosen product probability and representatives.



Writing \(B_t=I(e_t)\), the centered variables therefore satisfy \(\mathbb E[B_sB_t]=\min(s,t)\). The expectation is finite by Cauchy–Schwarz and the square integrability of the two variables.


\section{\texorpdfstring{Joint Gaussianity of the time family}{Joint Gaussianity of the time family}}


The isonormal isometry is linear and every output is scalar Gaussian. Apply Gaussianity of linear combinations of the isonormal map to the family \(t\mapsto e_t\). Every finite collection of the selected random variables is therefore jointly Gaussian.



Finite linear combinations of time values are isonormal evaluations at finite linear combinations of indicators. Their Gaussian laws establish the process's joint Gaussian property.


\section{\texorpdfstring{The Brownian finite-\allowbreak{}dimensional and increment properties}{The Brownian finite-dimensional and increment properties}}


A centered jointly Gaussian family with covariance \(\min(s,t)\) has Brownian finite-dimensional distributions. Differences over disjoint intervals have zero covariance, so their joint Gaussianity makes them independent. These are the increment and distribution properties required by the square-mean construction; sample-path continuity is not asserted here.



Centered Gaussianity and covariance \(\min(s,t)\) determine the Brownian finite-dimensional laws. Orthogonality of disjoint indicator increments gives independence of disjoint time increments.


\section{\texorpdfstring{The product sample space for spatial coordinates}{The product sample space for spatial coordinates}}


For a coordinate index set \(I\), take one scalar Brownian sample for each \(i\in I\). The product sample space is \(\Omega^I\), allowing independent entire coordinate processes rather than only independent values at one time.



The vector source is the finite product of scalar process sample spaces. Each spatial coordinate reads one separate factor, making the scalar paths independent by construction.


\section{\texorpdfstring{The vector source probability}{The vector source probability}}


Equip the vector sample space with the infinite product of the scalar Brownian source probabilities. The coordinate processes inherit independence from this product measure.



The finite product probability gives mass one and retains the scalar source law under each projection. It is the common probability space for the Euclidean Brownian vector.


\section{\texorpdfstring{A scalar spatial coordinate process}{A scalar spatial coordinate process}}


Define the \(i\)-th coordinate process by \(B_i(t,\omega)=B(t,\omega_i)\). It first projects to the \(i\)-th scalar sample, then evaluates the already constructed scalar process at \(t\).



A coordinate process is the scalar Brownian process composed with one product projection. This explicitly identifies which source factor supplies each spatial coordinate.


\section{\texorpdfstring{Brownian properties of each spatial coordinate}{Brownian properties of each spatial coordinate}}


The product coordinate projection has the scalar Brownian source as its law. Pull the scalar pre-Brownian process back along this projection; its finite-dimensional distributions, mean, and covariance are preserved. Thus each vector coordinate has the established scalar Brownian laws.



The coordinate projection preserves the scalar probability law. Hence all centered Gaussian increment and covariance properties transfer to the corresponding spatial coordinate.


\section{\texorpdfstring{Independence of the coordinate paths}{Independence of the coordinate paths}}


The product coordinate projections are independent as sample-valued maps. Compose each projection with the map sending a scalar sample to its entire time-indexed process. Independence is preserved under these coordinatewise functions, giving independence of the whole coordinate paths.



Independence is asserted for whole coordinate paths as observations of separate product factors. It therefore entails independence of every finite collection of observations drawn from distinct coordinates.


\section{\texorpdfstring{Independence after restriction to mesh times}{Independence after restriction to mesh times}}


Restrict each independent coordinate path to the deterministic grid \(nh\). This is a coordinatewise function of the path, so independence is preserved. The result supplies exactly the countable grid-path hypothesis used in the innovation theorems.



Restricting each path to mesh times is a measurable observation of that path. Independence is preserved by this restriction, so the grid paths retain spatial independence.


\section{\texorpdfstring{Normalization of the vector probability}{Normalization of the vector probability}}


Each scalar Brownian sample source has mass one. The infinite product over vector coordinates therefore also has mass one. This normalization concerns the vector product measure used to define the coordinate processes.



Normalization of the vector product source supplies the probability setting for the later Euclidean Gaussian and Brownian moment calculations. No additional weighting of the coordinate paths is introduced.


\end{document}
